\chapter{Michael Francis Atiyah}\label{ch1966michael-francis-atiyah:michael-francis-atiyah-fields-medal-1966-a-comprehensive-report}

\begin{center}\rule{0.5\linewidth}{0.5pt}\end{center}

\textbf{Prepared:} July 2026
\textbf{Subject:} Sir Michael Francis Atiyah (b. April 22, 1929, London, England; d. January 2019)
\textbf{Position:} Professor, Oxford University; Director, Isaac Newton Institute
\textbf{Doctoral Advisor:} William V. D. Hodge (Cambridge University)
\textbf{Fields Medal Citation:} \emph{``Did joint work with Hirzebruch in K-theory;
proved jointly with Singer the index theorem of elliptic operators on complex
manifolds; worked in collaboration with Bott to prove a fixed point theorem
related to the `Lefschetz formula'.''}

\begin{center}\rule{0.5\linewidth}{0.5pt}\end{center}

\begin{prereqs}
Algebraic topology (K-theory, characteristic classes), differential
geometry (vector bundles, connections), basic complex analysis.
\end{prereqs}

\begin{objectives}
State the Atiyah--Singer index theorem \cite{atiyah1963} and explain its significance;
describe K-theory as a generalized cohomology theory; explain the
connection between topology and physics.
\end{objectives}

\section{Problem History}

The index of a linear elliptic differential operator $D: C^\infty(E) \to
C^\infty(F)$ on a compact manifold $M$ is defined as:

\[\text{ind}(D) = \dim \ker D - \dim \text{coker} D\]

This is a homotopy invariant: operators with homotopic symbols have the same
index. The Russian school (Gelfand, Seeley, Duistermaat) established these
properties but no explicit formula was known.

The \textbf{Gauss--Bonnet theorem} expressed $\chi(M)$ as an integral of the
Pfaffian:
\[\chi(M) = \frac{1}{(2\pi)^{n/2}} \int_M \text{Pf}(R)\]
where $R$ is the curvature 2-form of a Riemannian metric.

The \textbf{Hirzebruch signature theorem} expressed $\tau(M)$ via the $L$-genus:
\[\tau(M) = \int_M L(TM)\]
where $L(TM) = \prod_i \frac{x_i}{\tanh x_i}$ and $x_i$ are the formal roots
of the curvature.

The \textbf{Hirzebruch--Riemann--Roch theorem} computed the Euler characteristic
of a holomorphic vector bundle $E$ on a complex manifold $M$:
\[\chi(M, E) = \int_M \text{ch}(E) \cdot \text{td}(M)\]
where $\text{td}$ is the Todd class.

These were special cases of a general formula connecting analysis to topology.
The bridge is the \textbf{symbol} of the operator: an elliptic operator $D$ has
a principal symbol $\sigma(D): \pi^*E \to \pi^*F$ on $T^*M$, which is an
isomorphism away from the zero section. This defines a class
$[\sigma(D)] \in K(T^*M)$.

K-theory had been introduced by Grothendieck in algebraic geometry and adapted
to topology by Atiyah and Hirzebruch. Bott periodicity gave it its
$\mathbb{Z}/2$-graded structure:
\[K(X \times \mathbb{R}^2) \cong K(X)\]
but the connection to elliptic operators was unclear.

The \textbf{Dirac operator} on a spin manifold $M$ satisfies
$\text{ind}(D) = \widehat{A}(M)$, known only in special cases. The
Lichnerowicz formula $D^2 = \nabla^*\nabla + \frac{1}{4}R$ showed positive
scalar curvature implies $\widehat{A}(M) = 0$.

The \textbf{Grothendieck--Riemann--Roch theorem} (with Hirzebruch) extended
the classical theorem to morphisms between varieties. For a proper map
$f: X \to Y$ between smooth varieties:
\[f_!(\text{ch}(E) \cdot \text{td}(X)) = \text{ch}(f_!E) \cdot \text{td}(Y)\]
This is the "formal" analogue of the index theorem in algebraic geometry.

The \textbf{Hirzebruch--Riemann--Roch theorem} for a vector bundle $E$ on a
complex manifold $M$ states:
\[\chi(M, E) = \sum_i (-1)^i \dim H^i(M, E) = \int_M \text{ch}(E) \cdot \text{td}(M)\]
The Todd class is given by $\text{td}(M) = \prod_i \frac{x_i}{1 - e^{-x_i}}$
where $x_i$ are the formal Chern roots of the tangent bundle.

The \textbf{Atiyah--Bott fixed point theorem} generalizes the Lefschetz formula
to elliptic complexes, but the general statement was not yet established.
The classical Lefschetz number for a map $f: M \to M$ is:
\[\Lambda(f) = \sum_i (-1)^i \text{Tr}(f^*|_{H^i(M;\mathbb{Q})})\]
The Lefschetz fixed point theorem states that $\Lambda(f)$ equals the sum
of local contributions at fixed points.

The broader question was the relationship between \emph{analysis} (kernel and
cokernel) and \emph{topology} (symbol class). The index is the bridge.

\section{How Was It Solved}

The \textbf{Atiyah--Singer index theorem} (1963, published 1968--1971) states:

\[\text{ind}(D) = \int_M \text{ch}(\sigma(D)) \cdot \widehat{A}(TM)\]

where $\sigma(D) \in K(T^*M)$ is the symbol class, $\text{ch}$ is the Chern
character, and $\widehat{A}(TM)$ is the $\widehat{A}$-genus. In K-theoretic
form, using the direct image map $p_!: K(T^*M) \to K(\text{pt}) = \mathbb{Z}$:

\[\text{ind}(D) = p_!([\sigma(D)])\]

\subsection*{The Proof Strategy}

The proof proceeds in three stages.

\textbf{Step 1: Reduction to Twisted Dirac Operators.}
Any elliptic operator $D$ with symbol $\sigma(D)$ can be written as
$D_0 \otimes 1_E$ where $D_0$ is the spin$^c$ Dirac operator and $E$ is a
vector bundle. By the Künneth theorem in K-theory, the symbol class of a
twisted Dirac operator is $[\sigma(D)] = [\sigma(D_0)] \cdot [E]$.

\textbf{Step 2: Heat Kernel Computation.}
For the Dirac operator, one computes the index using the heat equation.
Let $D$ be a self-adjoint Dirac operator. Then:
\[\text{ind}(D) = \text{Tr}(e^{-tD^2})|_{t \to 0}\]
The heat kernel $K(t, x, y)$ has an asymptotic expansion as $t \to 0$:
\[K(t, x, x) \sim (4\pi t)^{-n/2} \sum_{k=0}^\infty t^k a_k(x)\]
The leading term gives the index via the stationary phase approximation, and
the higher terms vanish in the supertrace.

\textbf{Step 3: Naturality and Multiplicativity.}
The index map $[\sigma(D)] \mapsto \text{ind}(D)$ is a homomorphism from
$K(T^*M)$ to $\mathbb{Z}$. By the naturality of the construction, it suffices
to verify the formula for a complete set of generators. The Thom isomorphism
and Bott periodicity provide these generators.

The key property is \textbf{multiplicativity}: if $M = M_1 \times M_2$, then:
\[\text{ind}(D_{M_1 \times M_2}) = \text{ind}(D_{M_1}) \cdot \text{ind}(D_{M_2})\]
This follows from the Künneth formula in K-theory and the product formula
for characteristic classes.

\subsection*{Special Cases}

The theorem recovers classical results:

\textbf{Gauss--Bonnet--Chern:} For $D = d + d^*$ on forms,
$\text{ind}(D) = \chi(M) = \int_M e(TM)$.

\textbf{Hirzebruch Signature:} For the signature operator,
$\text{ind}(D) = \tau(M) = \int_M L(TM)$.

\textbf{Hirzebruch--Riemann--Roch:} For the Dolbeault operator,
$\text{ind}(\bar{\partial}) = \int_M \text{ch}(E) \cdot \text{td}(M)$.

\textbf{Dirac Index:} For the Dirac operator on a spin manifold,
$\text{ind}(D) = \int_M \widehat{A}(TM)$.

The \textbf{Bott periodicity theorem} is fundamental to K-theory. It states:
\[K^0(S^2) \cong \mathbb{Z}, \quad K^1(S^2) = 0\]
The generator is the Bott element $\beta$, and the suspension isomorphism:
\[K^0(X) \cong K^1(X \times S^1) \cong K^0(X \times \mathbb{R}^2)\]
implies the $\mathbb{Z}/2$-graded periodicity of K-theory.

The Bott element can be constructed explicitly using the clutching function.
On $S^2 = \mathbb{C} \cup \{\infty\}$, the clutching function for the generator
is $z \mapsto z: S^1 \to U(1)$. Iterating this construction gives all K-theory
groups of spheres.

\subsection*{The Atiyah--Bott Fixed Point Theorem}

K-theory, developed with Hirzebruch, is the Grothendieck group $K(X)$ of
complex vector bundles over $X$. For compact $X$:
\[K(X) = \{\text{virtual bundles } [E] - [F]\}\]
with tensor product making it a ring. The reduced K-theory satisfies:
\[\widetilde{K}(S^2) \cong \mathbb{Z}\]
generated by the Bott element. Bott periodicity states:
\[K(X) \cong K(X \times \mathbb{R}^2)\]
which implies the $\mathbb{Z}/2$-grading:
\[K(X) = K^0(X) \oplus K^1(X)\]

\subsection*{The Atiyah--Hirzebruch Spectral Sequence}

To compute K-theory, Atiyah and Hirzebruch constructed a spectral sequence:
\[E_2^{p,q} = H^p(X; K^q(\text{pt})) \Rightarrow K^{p+q}(X)\]
This relates K-theory to ordinary cohomology, with differentials encoding
higher operations. For a CW-complex, this provides a systematic way to
compute K-groups from cell structure.

\subsection*{The Family Index Theorem}

For a family of elliptic operators parametrized by a compact base $B$, the
index lies in $K(B)$:
\[\text{ind}_B(D) \in K(B)\]
The theorem states that the Chern character of the family index is given by:
\[\text{ch}(\text{ind}_B(D)) = p_!(\text{ch}([\sigma(D)]) \cdot \widehat{A}(T^*M))\]
where $p: M \times B \to B$ is the projection. This was proved with Singer and
applied to moduli problems in gauge theory.

\subsection*{Characteristic Classes}

The proof requires precise formulas for characteristic classes. The
\textbf{Chern character} of a bundle $E$ with Chern roots $x_i$ is:
\[\text{ch}(E) = \sum_i e^{x_i}\]
The \textbf{Todd class} is:
\[\text{td}(E) = \prod_i \frac{x_i}{1 - e^{-x_i}}\]
The \textbf{$\widehat{A}$-genus} for a spin manifold with Pontryagin roots $p_i$ is:
\[\widehat{A}(M) = \prod_i \frac{\sqrt{p_i/4}}{\sinh(\sqrt{p_i/4})}\]

These classes satisfy multiplicative properties essential to the proof.
The $\widehat{A}$-genus is a rational characteristic class, and for spin
manifolds it is integral. This integrality is a deep result related to
the index theorem itself.

\subsection*{The $\eta$-Invariant and APS Boundary Theorem}

Atiyah, Patodi, and Singer extended the index theorem to manifolds with
boundary. For a manifold $M$ with boundary $\partial M$, the index formula
acquires a boundary correction:
\[\text{ind}(D) = \int_M \widehat{A}(TM) - \frac{1}{2}(\eta(0) + h)\]
where $\eta(s) = \sum_\lambda \text{sign}(\lambda)|\lambda|^{-s}$ is the eta
invariant of the boundary operator $B$, and $h = \dim \ker B$.

The $\eta$-invariant measures the spectral asymmetry of the boundary operator
and is a subtle invariant appearing in many areas of geometry and physics.
It is related to the spectral flow and plays a key role in the study of
metrics with positive scalar curvature.

\subsection*{ADHM Construction}

With Drinfeld, Hitchin, and Manin, Atiyah gave a complete description of
instantons on $\mathbb{R}^4$ via quaternionic linear algebra. Instantons of
charge $k$ are parametrized by the moduli space $\mathcal{M}_k$ of dimension
$8k - 3$, constructed from solutions to the ADHM equations. This was crucial
for physics applications of the index theorem.

The ADHM construction identifies $\mathcal{M}_k$ with the quotient of a
variety in $\text{Mat}_{k+1,k}(\mathbb{H}) \times \text{Mat}_{k,k+1}(\mathbb{H})$
by the action of $U(k)$. The moduli space is smooth for $k \geq 1$ and has
the structure of a hyperkähler manifold. The ADHM construction has
since been generalized to other gauge groups and to instantons on
hyperkähler 4-manifolds.

\section{Impact of This Discovery}

The index theorem is one of the most important results of 20th-century
mathematics, unifying analysis, topology, and geometry.

\textbf{Gauge Theory:} The index theorem computes dimensions of instanton and
monopole moduli spaces. For a principal $G$-bundle $P$ over a 4-manifold, the
moduli space $\mathcal{M}$ of anti-self-dual connections has dimension given by
the index of the deformation complex. This underlies Donaldson's theorem on
smooth 4-manifolds and the Seiberg--Witten invariants.

\textbf{Anomalies in Physics:} Quantum field theories can have classical
symmetries broken by quantum effects. The index theorem explains gauge anomalies:
a gauge anomaly cancels iff a certain family index vanishes. This led to the
understanding of chiral anomalies in the Standard Model.

\textbf{Mathematical Physics:} The index theorem connects to the heat kernel
expansion, providing a bridge between physics and pure mathematics. The
Atiyah--Singer theorem explains the conformal anomaly in even-dimensional
quantum field theories.

\textbf{Number Theory:} The theorem connects to the Kronecker limit formula and
automorphic forms. The arithmetic applications arose through the study of
modular forms and $L$-functions.

\textbf{K-theory:} Topological K-theory became a major tool in algebraic
topology, algebraic geometry, and operator algebras. Kasparov's KK-theory is
fundamental to the Baum--Connes conjecture.

\textbf{Representation Theory:} The Atiyah--Bott fixed point theorem has deep
applications to representation theory of Lie groups and symplectic geometry.

\textbf{TQFT:} Atiyah's axiomatization of topological quantum field theories as
functors from cobordisms to vector spaces underlies the Jones polynomial and
Chern--Simons theory. He defined a TQFT as a symmetric monoidal functor:
\[\mathcal{Z}: \text{Cob}_n \to \text{Vect}\]
assigning vector spaces to $(n-1)$-manifolds and linear maps to $n$-manifolds.
The Jones polynomial arises from a 3-dimensional TQFT.

\textbf{Index Theorems for Manifolds with Corners:}
Atiyah, Patodi, and Singer developed boundary conditions for Dirac operators
on manifolds with corners. The local Atiyah--Patodi--Singer boundary condition
uses spectral projections of the boundary operator.

\textbf{Transversality and Geometry:}
The index theorem has deep connections to transversality theory. The
Thom--Porteous formula expresses degeneracy loci of bundle maps as
Schur polynomials in Chern classes, providing geometric interpretations
of index formulas.

\section{Modern Views and Improvements}

The index theorem extends to orbifolds, foliations (Connes--Skandalis), and
noncommutative geometry (Connes).

\textbf{Equivariant K-theory:} For a compact Lie group $G$ acting on $X$, the
equivariant K-theory $K_G(X)$ carries a module structure over the representation
ring $R(G)$. Bott periodicity extends equivariantly:
\[K_G^0(X) \cong K_G^0(X \times \mathbb{R}^2)\]
The equivariant index theorem relates the character of the representation on
the kernel to the fixed point data.

\textbf{KR-theory:} Atiyah introduced KR-theory for spaces with involution,
satisfying $\mathbb{Z}/8$-periodicity. This is important for real K-theory and
physical applications involving time-reversal symmetry.

\textbf{Baum--Connes Conjecture:} This conjecture relates the K-theory of the
C*-algebra of a discrete group $\Gamma$ to the K-homology of the classifying
space $E\Gamma$. A proof would imply the Novikov conjecture on homotopy
invariance of higher signatures and the Gromov--Lawson conjecture on positive
scalar curvature.

\textbf{Positive Scalar Curvature:} Studies use the Lichnerowicz obstruction
and Rosenberg index, guided by the index theorem. The $\widehat{A}$-genus
remains the fundamental obstruction on spin manifolds.

\textbf{Noncommutative Geometry:} Connes developed the index theorem in the
context of noncommutative manifolds, where the "space" is replaced by a
noncommutative algebra. The formula involves cyclic cohomology and the Chern
character in noncommutative geometry.

\textbf{TQFT and Chern--Simons:} The index-theoretic foundations developed by
Atiyah led to the rigorous construction of Chern--Simons theory and the
Jones polynomial. The Seiberg--Witten invariants, defined via monopole equations,
also rest on index-theoretic analysis.

\textbf{Dirac Operators on Singular Spaces:} Extensions to stratified spaces,
orbifolds, and spaces with conical singularities have been developed using
weighted Sobolev spaces and index theory on manifolds with edges.

\textbf{Arithmetic Applications:} The arithmetic Atiyah--Singer theorem relates
analytic torsion to Reidemeister torsion, with applications to L-functions
and the Stark conjectures.

\textbf{Noncommutative Geometry:} Connes further developed the index theorem
in noncommutative settings. For a noncommutative space defined by an algebra
$A$, the index pairing between $K_0(A)$ and $K^0(A)$ is computed by a
trace formula involving the cyclic cocycle associated to the Dirac operator.

\textbf{Physical Applications:} The index theorem explains anomalies in
quantum field theory. The chiral anomaly in 4 dimensions is computed by
the integral of the Chern character of the gauge field strength, giving
the famous result $\partial_\mu j^\mu_5 = \frac{e^2}{16\pi^2} F \wedge F$.

\begin{center}\rule{0.5\linewidth}{0.5pt}\end{center}

\section*{Summary}
\addcontentsline{toc}{section}{Summary}

Atiyah proved the Atiyah--Singer index theorem, connecting the analytical index
of elliptic operators to topological invariants of the underlying manifold.
Together with Hirzebruch, he developed topological K-theory, and with Bott he
proved the generalized Lefschetz fixed point theorem, unifying analysis,
topology, and geometry. The family index theorem and ADHM construction extended
these ideas to moduli problems and gauge theory, with profound implications for
mathematical physics.

The index theorem has become a fundamental tool connecting diverse areas of
mathematics, from differential geometry to algebraic topology, from number
theory to quantum field theory. Its reach continues to grow as new
generalizations and applications are discovered.
